% distributed-runs.tex — every cut plans the run before it: a chain that folds twice expands and contracts twice, each
% run shipped through its own branch and recombined at home, where the next run starts. Rendered by
% `bin/render-tikz.py --png --dpi 300 docs/skills/mtron/assets/tikz`, which wraps this body in a standalone document;
% the PNG goes to docs/website/skills/mtron/assets/. INK: white on a transparent background (the site theme is dark);
% the lanes are GREEN; a fold is a circle.

\begin{tikzpicture}[
  inst/.style={draw, rounded corners=2pt, thick, minimum width=11mm, minimum height=8mm, font=\small, color=white},
  fold/.style={draw, circle, thick, minimum size=9mm, font=\small, color=white},
  hop/.style={thick, color=white, -{Stealth[length=2mm]}},
  lane/.style={thick, color=green!60!white, -{Stealth[length=1.8mm]}},
  note/.style={font=\scriptsize, color=white},
  link/.style={draw=green!60!white, fill=green!18!black, thick, rounded corners=0.5pt, minimum width=2.8mm,
               minimum height=1.9mm, inner sep=0pt}
]

% ---- the chain up to the first cut ----
\node[inst] (f1) at (0,0)   {$f_1$};
\node[inst] (f2) at (2.2,0) {$f_2$};
\node[inst] (f3) at (4.4,0) {$f_3$};
\draw[hop] (f1) -- (f2);
\draw[hop] (f2) -- (f3);

% ---- the first run ships, and its fold recombines at home ----
\node[fold] (r1) at (9.0,0) {$r_1$};
\foreach \ang in {-60,-40,-20,0,20,40,60}
  \draw[lane] (f3.east) to[bend left=\ang] (r1.west);
% a few of the lane's own instructions, scattered along the lanes: each lane carries the run, a chain
\path (f3.east) to[bend left=60] node[link, pos=0.38, sloped] {} node[link, pos=0.62, sloped] {} (r1.west);
\path (f3.east) to[bend left=20] node[link, pos=0.5, sloped] {} (r1.west);
\path (f3.east) to[bend left=-20] node[link, pos=0.35, sloped] {} node[link, pos=0.65, sloped] {} (r1.west);
\path (f3.east) to[bend left=-60] node[link, pos=0.5, sloped] {} (r1.west);

% ---- home again: the next run starts where the combiner left it ----
\node[inst] (f4) at (11.2,0) {$f_4$};
\draw[hop] (r1) -- (f4);

% ---- the second run ships, and its fold recombines at home ----
\node[fold] (r2) at (15.8,0) {$r_2$};
\foreach \ang in {-60,-40,-20,0,20,40,60}
  \draw[lane] (f4.east) to[bend left=\ang] (r2.west);
% a few of the lane's own instructions, scattered along the lanes: each lane carries the run, a chain
\path (f4.east) to[bend left=60] node[link, pos=0.5, sloped] {} (r2.west);
\path (f4.east) to[bend left=0] node[link, pos=0.35, sloped] {} node[link, pos=0.65, sloped] {} (r2.west);
\path (f4.east) to[bend left=-40] node[link, pos=0.42, sloped] {} (r2.west);
\path (f4.east) to[bend left=40] node[link, pos=0.6, sloped] {} (r2.west);

% ---- the tail stays home ----
\node[inst] (f5) at (18.0,0) {$f_5$};
\draw[hop] (r2) -- (f5);

% ---- expansion and contraction, twice ----
\node[note] at (6.7,-1.45)  {expand};
\node[note] at (9.0,-1.45)  {contract};
\node[note] at (13.5,-1.45) {expand};
\node[note] at (15.8,-1.45) {contract};

\end{tikzpicture}
